\documentclass[11pt,a4paper]{article}
\usepackage[utf8]{inputenc}
\usepackage[T1]{fontenc}
\usepackage{geometry}
\usepackage{xcolor}
\usepackage{hyperref}
\usepackage{titlesec}

\geometry{margin=1in}
\definecolor{tumblue}{RGB}{0,101,189}

\hypersetup{
    colorlinks=true,
    linkcolor=tumblue,
    urlcolor=tumblue,
}

\titlelabel{\thetitle.\quad}

\begin{document}

\begin{center}
    {\LARGE\bfseries\color{tumblue} Research Proposal Outline}\\[6pt]
    {\large Robust Spatial-Temporal Fusion: Graph-State Space Models (G-SSM)}\\[4pt]
    \textbf{Ismail AISSAOUI} $|$ State Engineer in Artificial Intelligence
\end{center}

\bigskip

\section{Introduction}
Modern machine learning for complex systems (Climate, Industry, Finance) requires models that can capture both intricate spatial dependencies and long-range temporal dynamics. While Transformers and standard GNNs are the current baseline, they face significant scaling and robustness issues. I propose to investigate a novel fusion of **Graph Neural Networks (GNNs)** and **State Space Models (SSMs)**, creating a **"Graph-State Space Model" (G-SSM)** framework designed for robust, multi-scale spatial-temporal forecasting.

\section{Proposed Research Axes}

\subsection{1. Graph-Integrated State Space Transitions}
Existing SSMs (like Mamba) operate on 1D sequences. I propose to generalize the SSM state transition matrix $A$ to be a function of the graph topology. By integrating graph-convolutional operators into the state space recurrence, we can capture how latent states "flow" across a network (e.g., thermal propagation in a gas turbine or weather patterns across a grid).

\subsection{2. Robustness to Noisy and Incomplete Graphs}
A major challenge in GNNs is sensitivity to edge noise or missing data. I intend to explore how the inherent stability properties of linear state space systems can be used to regularize graph-based embeddings. By treating the graph-state evolution as a dynamical system, we can apply control-theoretic stability constraints to ensure robust inference even when the underlying graph is partially corrupted.

\subsection{3. Probabilistic G-SSMs for Uncertainty Quantification}
For high-stakes applications, point estimates are insufficient. I propose to develop probabilistic versions of the G-SSM that provide calibrated **Uncertainty Quantification (UQ)**. Building on my work with hybrid modeling at Sonatrach, I will investigate how domain-priors can be used to bound the uncertainty of the graph-state transitions, particularly in the "tails" of the distribution where extreme events occur.

\section{Alignment with TUM CIT}
This research aligns with the Data Analytics and ML group's focus on robust and theoretically sound graph learning. I aim to validate these G-SSM architectures on large-scale datasets (Sensor networks, Climate grids), contributing to the development of next-generation foundational models for complex spatial-temporal systems.

\end{document}
